\documentclass[10pt,a4paper]{amsart}
\usepackage[margin=19mm]{geometry}
\usepackage{amsmath,amssymb,mathtools}
\usepackage[scr=rsfs,cal=euler]{mathalfa}
\usepackage{xfrac}
\usepackage[unicode,hidelinks,bookmarks=false]{hyperref}
\newcommand{\ie}{i.e.\ }
\newcommand{\cf}{cf.\ }
\newcommand{\resp}{respectively\ }
\newcommand{\Subsection}{\subsection}
\providecommand{\nobreakdash}{\nobreak\hbox{-}}
\setlength{\emergencystretch}{2em}
\multlinegap0pt
\usepackage{amssymb}
\usepackage[shortlabels]{enumitem}
\newtheorem{lemma}{Lemma}
\title{On the discrete mean square of certain hybrid sum involving $a_{\mathbb{K}}(n)$}
\author{EKTA SONI, M.S. DATT, AYYADURAI SANKARANARAYANAN}
\date{Source: arXiv:2603.11621v1 (CC BY 4.0)}
\begin{document}
\maketitle

\noindent\emph{Source and licence.} On the discrete mean square of certain hybrid sum involving $a_{\mathbb{K}}(n)$. Version arXiv:2603.11621v1 (CC BY 4.0), distributed by arXiv under the Creative Commons Attribution 4.0 International licence, http://creativecommons.org/licenses/by/4.0/, as recorded in the arXiv licence metadata for this exact version and shown on the abstract page https://arxiv.org/abs/2603.11621v1. Full licence notice: \texttt{licenses/arxiv-2603.11621-CC-BY-4.0.txt}.\par\smallskip

\noindent\textit{Editorial scope.} Counted unit: the article's lemma lemma (source0.tex lines 265--267). The article's complete original proof of it is delivered with it (source0.tex lines 268--318, one \texttt{proof} environment of 51 lines). No mathematics is written, repaired or re-derived: the statement and the proof are the article's own lines, in the article's own notation, and no retained line is substituted or re-flowed.  No further source-local context is needed: every cross-reference of every retained line resolves inside the two delivered spans.  Nothing was omitted from any retained span, so the package carries no omission record.  No retained line contains a citation, so the wrapper carries no bibliography and no bibliography entry.  No retained line contains a figure environment, an image-inclusion command or drawing code, so no image is delivered and the package declares no asset dependency.  Editorial formatting only, in addition: an AMS \texttt{amsart} preamble that restates the article's own declarations and macros used by the retained lines, namely the environments \texttt{lemma} on separate counters, together with the standard packages those lines need.  The retained environments print in this document as Lemma 1 in order of appearance, whereas the article prints the same items with the numbering of its own full document.  The complete native source of the article as distributed in the arXiv e-print for this exact version is bundled unchanged as \texttt{sources/196227/source0.tex}.
\begin{lemma}
	Let $g(n)=(-1)^{n}\sum_{d|n}(-1)^{d}d^{3}$. Then $g(n)$ is a multiplicative arithmetic function.
\end{lemma}

\begin{proof}
	To prove the lemma, we show that for any two distinct primes $p,\,q$ and for $\ell,m \in \mathbb{N}$
	\begin{center}
		$g(p^{\ell}q^{m})=g(p^{\ell})g(q^{m})$.
	\end{center}
	\textbf{Case(i)} In this case we take $p=2$ and $q$ is odd. By the definition of $g(n)$,
	\begin{align*}
		g(2^{\ell})&=(-1)^{2^{\ell}}\{-1+2^{3}+2^{6}+\dots+2^{3\ell}\}\\
		&=\{-1+2^{3}+2^{6}+\dots+2^{3\ell}\},
	\end{align*}
	and 
	\begin{align*}
		g(q^{m})&=(-1)^{q^{m}}\{-1-q^{3}-\dots-q^{3m}\}\\
		&=\{1+q^{3}+\dots+q^{3m}\}.
	\end{align*}
	Therefore,
	\begin{align*}
		g(2^{\ell})g(q^{m})&=\{-1+2^{3}+2^{6}+\dots+2^{3\ell}\}\{1+q^{3}+\dots+q^{3m}\}\\
		&=-\{1+q^{3}+\dots+q^{3m}\}+\sum\limits_{\substack{1\leq i\leq \ell \\ 0\leq j \leq m}} 2^{3i} q^{3j}\\
		&=-\sum\limits_{j=0}^{m}q^{3j}+\sum\limits_{\substack{1\leq i\leq \ell \\ 0\leq j \leq m}} 2^{3i} q^{3j}.
	\end{align*}
	Now take,
	\begin{align*}
		g(2^{\ell}q^{m})&=(-1)^{2^{\ell}q^{m}}\big\{-1+\sum\limits_{i=1}^{\ell}2^{3i}-\sum\limits_{j=1}^{m}q^{3j}+\sum\limits_{\substack{1\leq i\leq \ell \\ 1\leq j \leq m}} 2^{3i} q^{3j}\big\}\\
		&=-\big\{1+\sum\limits_{j=1}^{m}q^{3j}\big\}+\big\{\sum\limits_{i=1}^{\ell}2^{3i}+ \sum\limits_{\substack{1\leq i\leq \ell \\ 1\leq j \leq m}} 2^{3i} q^{3j}\big\}\\
		&=-\sum\limits_{j=0}^{m}q^{3j}+\sum\limits_{\substack{1\leq i\leq \ell \\ 0\leq j \leq m}} 2^{3i} q^{3j}.
	\end{align*}
	Therefore, $g(2^{\ell}q^{m})=g(2^{\ell})g(q^{m})$.\vspace{3mm}\\
	\textbf{Case (ii)} Suppose both $p,\,q$ are odd primes.\\
	We have,
	\begin{align*}
		g(p^{\ell})&=(-1)^{p^{\ell}}\{-1-p^{3}-p^{6}\dots-p^{3\ell}\}\\
		&=1+p^{3}+p^{6}+\dots+p^{3\ell},
	\end{align*}
	and 
	\begin{align*}
		g(p^{m})&=1+q^{3}+q^{6}+\dots+q^{3m}.
	\end{align*}
	Therefore,
	\begin{align*}
		g(p^{\ell})g(p^{m})&=\{1+p^{3}+p^{6}+\dots+p^{3\ell}\}\{1+q^{3}+q^{6}+\dots+q^{3m}\}\\
		&=\sum\limits_{\substack{0\leq i\leq \ell \\ 0\leq j \leq m}} p^{3i}q^{3j}.
	\end{align*}
	Now, 
	\begin{align*}
		g(p^{\ell}q^{m})=(-1)^{p^{\ell}q^{m}}\sum\limits_{\substack{0\leq i\leq \ell \\ 0\leq j \leq m}}(-1)^{p^{3i}q^{3j}} p^{3i}q^{3j}=\sum\limits_{\substack{0\leq i\leq \ell \\ 0\leq j \leq m}} p^{3i}q^{3j}.
	\end{align*}
	Therefore,
	$g(p^{\ell}q^{m})=g(p^{\ell})g(q^{m})$.\\
	Hence the Lemma.
\end{proof}

\end{document}
